\documentclass[conference]{IEEEtran}
\usepackage{cite}
\usepackage{amsmath,amssymb,amsfonts}
\usepackage{algorithmic}
\usepackage{graphicx}
\usepackage{textcomp}
\usepackage{xcolor}
\usepackage{enumitem}
\usepackage{flushend}
\usepackage{bm}

\def\BibTeX{{\rm B\kern-.05em{\sc i\kern-.025em b}\kern-.08em
    T\kern-.1667em\lower.7ex\hbox{E}\kern-.125emX}}

\ifCLASSOPTIONcompsoc
\usepackage[caption=false,font=normalsize,labelfont=sf,textfont=sf]{subfig}
\else
\usepackage[caption=false,font=footnotesize]{subfig}
\fi

\DeclareMathOperator{\f}{\textit{f}}

\begin{document}

\title{
Development and Analysis of a DCT-Based Flipbook Compression Pipeline with Quantization and Huffman Coding
}

\author{
    \IEEEauthorblockN{Artem Onyshchuk, Yarema Mykhasiak, Yulian Volodymyr Zaiats}
    \IEEEauthorblockA{\textit{Faculty of Applied Sciences} \\
    \textit{Ukrainian Catholic University}, Lviv, Ukraine  \\
    \{a.onyshchuk.pn, mykhasiak.pn, zaiats.pn\}@ucu.edu.ua}
    \and
    \IEEEauthorblockN{Viktor Syrotiuk}
    \IEEEauthorblockA{\textit{Faculty of Applied Sciences} \\
    \textit{Ukrainian Catholic University}\\
    syrotiuk.pn@ucu.edu.ua}
    \and
    \IEEEauthorblockN{Denys Maletskyy}
    \IEEEauthorblockA{\textit{Graphics Engineer, Zibra AI} \\
    Lviv, Ukraine \\
    maletsden@gmail.com}
}

\maketitle

\begin{abstract}
Efficient flipbook compression is essential for high-frequency visual graphics where memory bandwidth is a limiting factor. This paper documents the evolutionary development of a high-performance flipbook compression pipeline. We transition from an initial hybrid CPU-GPU architecture to a fully GPU-accelerated system that offloads both transform coding and entropy coding to NVIDIA CUDA. By utilizing a 3-stage parallel bit-packing algorithm, we achieve a compression ratio of 11.2x and a throughput of 490+ FPS for compression and 2400+ FPS for decompression at 1080p. Our analysis includes a rigorous evaluation of memory bandwidth, occupancy, and asynchronous scheduling, demonstrating the scalability of the architecture.
\end{abstract}

\begin{IEEEkeywords}
HPC, CUDA, Image Compression, Discrete Cosine Transform, Parallel Huffman, Occupancy Analysis, Bandwidth.
\end{IEEEkeywords}

\section{Introduction}

Digital flipbooks are sequences of static images that create animation through rapid playback. Traditional RAW storage methods incur prohibitive memory costs. While video compression (inter-frame) achieves high ratios, intra-frame compression is often preferred for flipbooks to ensure low-latency random access to individual animation states. This project develops a DCT-based compression pipeline adhering to the HPC:NE performance standards ($>10$x ratio, $>20$ dB PSNR, $>60$ FPS).

\section{Mathematical Foundations}

\subsection{Luminance and Chrominance Transformation}
Frames are converted to the YCbCr space using the ITU-R BT.601 standard:
\begin{equation}
\begin{bmatrix} Y \\ C_b \\ C_r \end{bmatrix} = 
\begin{bmatrix} 
0.299 & 0.587 & 0.114 \\
-0.1687 & -0.3313 & 0.5 \\
0.5 & -0.4187 & -0.0813
\end{bmatrix}
\begin{bmatrix} R \\ G \\ B \end{bmatrix}.
\end{equation}
Chroma planes are subsampled at a $4{:}2{:}0$ ratio to reduce spatial redundancy.

\subsection{Discrete Cosine Transform (DCT)}
Image planes are partitioned into $N \times N$ blocks. The 2D DCT-II transform is defined as:
\begin{equation}
C(u,v) = \frac{2}{N}\alpha(u)\alpha(v) \sum_{x=0}^{N-1}\sum_{y=0}^{N-1} X'(x,y) \cos \Phi(x,u) \cos \Phi(y,v),
\end{equation}
where $\Phi(i, j) = \frac{(2i+1)j\pi}{2N}$. This transform effectively concentrates spectral energy in a small subset of coefficients. While Fast Fourier Transform (FFT) was considered for frequency decomposition, DCT-II was selected due to its superior energy compaction properties for natural image signals and its real-valued coefficients which simplify parallel implementation on CUDA. 

\subsection{Quantization and Zigzag Scanning}
Coefficients $C(u,v)$ are quantized using psycho-visual matrices. To maximize the efficiency of entropy coding, blocks are converted into a 1D vector via a \textbf{zigzag scan}, grouping zeros at the end of the sequence.

\section{Evolution of the Pipeline Architecture}

The development of the codec progressed through two distinct phases, reflecting an evolution from basic parallelization to full hardware offloading.

\subsection{Phase I: Initial Hybrid Implementation}
In the first phase, a hybrid architecture was adopted. Dense numerical operations (Color Conversion, DCT, and Quantization) were executed as GPU kernels. However, entropy coding (RLE and Huffman) remained on the CPU.
\begin{itemize}[leftmargin=*]
    \item \textbf{Rationale}: Entropy coding involves irregular branching and sequential bit-packing, which are theoretically difficult to parallelize on SIMT hardware.
    \item \textbf{Limitation}: The resulting pipeline was bottlenecked by frequent Device-to-Host (D2H) transfers of quantized coefficients. The PCIe bus overhead limited total throughput to approximately 32-45 FPS.
\end{itemize}

\subsection{Phase II: Full GPU Acceleration (Final Architecture)}
To overcome the synchronization wall, we implemented a \textbf{Full GPU Pipeline}. This phase introduced the parallel entropy coding backend:
\begin{itemize}[leftmargin=*]
    \item \textbf{Parallel Bit-packing}: Uses a 3-kernel strategy (Analysis $\to$ Scan $\to$ Pack) to resolve symbol offsets in parallel.
    \item \textbf{Result}: By eliminating D2H synchronization for intermediate data, compression throughput reached \textbf{490+ FPS}, a 10x improvement over Phase I.
\end{itemize}

\section{GPU Resource and Optimization Analysis}

\subsection{Occupancy and Shared Memory}
For $32 \times 32$ tiles, the system utilizes \textbf{12 KiB} of Shared Memory for block staging. By keeping register pressure under 32 per thread, we achieved an occupancy of \textbf{74.8\%} on Ampere hardware.

\subsection{Memory Bandwidth Utilization}
The DCT kernel achieves a compute bandwidth of \textbf{9.67 GB/s}. The parallel packing stage reaches \textbf{2.01 GB/s}. Constant Memory is used for DCT kernels and quantization tables, ensuring high cache hit rates for frequently coefficient lookups.

\subsection{Asynchronous Execution Model}
We implemented an asynchronous execution model using pinned memory and dual CUDA streams. This overlaps the transfer of frame $k+1$ with the processing of frame $k$, hiding \textbf{92.4\%} of the PCIe communication latency.

\section{Experimental Results}

We evaluated the performance of the three implementations (Serial, OpenMP, and CUDA).

\begin{figure}[htbp]
    \centering
    \includegraphics[width=0.48\textwidth]{chart_summary_table.png}
    \caption{Comparative Performance and Quality Summary. Shows the dramatic throughput increase from Phase I (Hybrid) to Phase II (Full GPU).}
    \label{fig:summary_table}
\end{figure}

The rate-distortion characteristics (Fig. \ref{fig:psnr_quality} and \ref{fig:rate_distortion}) demonstrate a compression ratio of \textbf{11.2x} at a stable PSNR of 37-54 dB. We also evaluated the Structural Similarity Index (SSIM), which remained consistently above 0.98 for $Q > 50$, indicating excellent perceptual preservation.

\begin{figure}[htbp]
    \centering
    \includegraphics[width=0.45\textwidth]{chart_psnr_quality.png}
    \caption{PSNR vs. Quality Tier. Stability is consistent across all test sets.}
    \label{fig:psnr_quality}
\end{figure}

\begin{figure}[htbp]
    \centering
    \includegraphics[width=0.45\textwidth]{chart_rate_distortion.png}
    \caption{Rate-Distortion curve for the final full-GPU pipeline.}
    \label{fig:rate_distortion}
\end{figure}

The final speedup reached \textbf{35.8x} compared to the serial baseline.

\begin{figure}[htbp]
    \centering
    \includegraphics[width=0.45\textwidth]{chart_speedup.png}
    \caption{Speedup factor comparison. Shows the evolutionary jump in performance through CUDA parallelism.}
    \label{fig:speedup}
\end{figure}

\section{Conclusion}

This project documents the successful evolution of a flipbook compression pipeline. By transitioning from a hybrid CPU-GPU model to a fully offloaded parallel architecture, we eliminated the entropy coding bottleneck. The final system delivers deterministic 2400+ FPS decompression, easily meeting the high-performance requirements of modern visual systems.

\begin{thebibliography}{00}
\bibitem{b1} G.~K. Wallace, ``The JPEG still picture compression standard,'' \emph{Communications of the ACM}, 1991.
\bibitem{b2} NVIDIA, ``CUDA C++ Programming Guide,'' v12.0, 2023.
\bibitem{b3} D.~A. Huffman, ``A method for the construction of minimum-redundancy codes,'' \emph{Proceedings of the IRE}, 1952.
\bibitem{b4} J.~Sanders, \emph{CUDA by Example: An Introduction to General-Purpose GPU Programming}, 2010.
\bibitem{b5} S.~Fujita, ``TinyEXR library for OpenEXR file I/O,'' GitHub, 2023.
\bibitem{b6} S.~Barrett, ``stb\_image: Single-file libraries for C/C++,'' GitHub, 2023.
\bibitem{b7} NVIDIA CUB Library, ``CUDA Unbound (CUB) for Scan and Reduction primitives,'' GitHub, 2023.
\bibitem{b8} K.~R. Rao, \emph{Discrete Cosine Transform: Algorithms, Advantages, Applications}, Academic Press, 1990.
\bibitem{b9} W.~K. Pratt, \emph{Digital Image Processing}, 4th~ed. Wiley-Interscience, 2007.
\bibitem{b10} V.~Bhaskaran and K.~Konstantinides, \emph{Image and Video Compression Standards: Algorithms and Architectures}, 1997.
\end{thebibliography}

\end{document}
